\documentclass[11pt]{article}

\usepackage[utf8]{inputenc}
\usepackage[T1]{fontenc}
\usepackage{geometry}
\usepackage{amsmath,amssymb,amsfonts}
\usepackage{booktabs}
\usepackage{graphicx}
\usepackage{subcaption}
\usepackage{multirow}
\usepackage{hyperref}
\usepackage{enumitem}
\usepackage{xcolor}
\usepackage{array}
\usepackage{mathtools}

\geometry{margin=1in}
\hypersetup{
  colorlinks=true,
  linkcolor=blue,
  citecolor=blue,
  urlcolor=blue
}

\title{SCP-Flow: Structure-Controllable Progression Flow for Longitudinal Glaucoma Forecasting}
\author{Draft}
\date{\today}

\begin{document}

\maketitle

\begin{abstract}
Glaucoma forecasting from longitudinal fundus imaging requires a future prediction model that preserves temporal progression patterns, leverages population-level priors, and remains robust under missing clinical follow-up. This paper framework introduces Structure-Controllable Progression Flow (SCP-Flow), a flow-based extension of tHPM-LDM that replaces latent diffusion future prediction with a history-conditioned progression flow while preserving the original historical and population conditioning modules. The framework is designed to support latent image forecasting, structure-aware conditioning, next-visit interval prediction, uncertainty estimation, and missing-visit robustness analysis.
\end{abstract}

\tableofcontents

\section{Introduction}

Glaucoma forecasting from longitudinal fundus imaging requires a model that can preserve historical trajectory information while remaining robust to irregular follow-up and missing visits. The current codebase starts from tHPM-LDM and upgrades its future prediction backbone from latent diffusion to a progression flow model.

This draft documents the current SCP-Flow implementation and serves as the starting point for a paper-formatted manuscript.

\section{Related Work}

\subsection{Longitudinal Glaucoma Forecasting}

Longitudinal glaucoma forecasting aims to predict future visual status, disease progression, or future retinal appearance from historical ophthalmic observations. Prior work has explored both image-based and clinical time-series based approaches, with increasing emphasis on irregular visit intervals and patient-specific progression patterns.

\subsection{Latent Generative Forecasting}

Recent work has shown that latent-space generative models provide an effective way to model future ophthalmic images while reducing the computational burden of pixel-space generation. Diffusion-based forecasting models can capture complex future uncertainty, but often require expensive sampling and may not expose a direct progression field.

\subsection{Flow-Based Temporal Modeling}

Flow-based trajectory models directly parameterize temporal evolution and can provide a more explicit representation of disease progression in latent space. In the glaucoma setting, this is attractive because the disease often evolves gradually over time, making progression-aware latent transport a natural modeling choice.

\subsection{Structure-Aware Retinal Modeling}

Structural biomarkers such as cup-to-disc ratio and optic disc morphology are central to glaucoma assessment. ROI-based conditioning and structure-guided generation therefore offer a natural path toward more interpretable longitudinal forecasting.

\section{Methods}

\subsection{Problem Setup}

Given a longitudinal SIGF clip
\begin{equation}
\mathcal{S}=\{(x_1,l_1,t_1),\dots,(x_T,l_T,t_T)\},
\end{equation}
where $x_i$ is the $i$-th visit image, $l_i$ is the visit label, and $t_i$ is the visit time, the goal is to predict the future visit state using the preceding history.

In the current implementation, the clip length is fixed to
\begin{equation}
T=6.
\end{equation}

\subsection{Latent Representation}

We reuse the pretrained first-stage autoencoder from tHPM-LDM. Each visit image is mapped into the latent space by
\begin{equation}
z_i = E(x_i),
\end{equation}
and the predicted future latent is decoded by
\begin{equation}
\hat{x}_T = D(\hat{z}_T).
\end{equation}

\subsection{Historical and Structural Conditioning}

The current SCP-Flow implementation keeps the original historical and population conditions:
\begin{equation}
(c_h, c_p) = \mathrm{CondGen}(x_{1:T-1}, t_{1:T-1}, l_{1:T-1}).
\end{equation}

In addition, the model uses structure-related conditions from the last observed visit, including scalar structure descriptors and ROI crops:
\begin{equation}
s_{\mathrm{scalar}} = [\mathrm{vCDR}, \mathrm{OCOD}],
\end{equation}
\begin{equation}
c_s = \mathrm{StructEnc}(s_{\mathrm{scalar}}, r_{\mathrm{disc}}, r_{\mathrm{polar}}).
\end{equation}

\subsection{Bridge Sampling and Flow Prediction}

During training, an intermediate latent is sampled between the previous visit latent $z_{T-1}$ and the future visit latent $z_T$:
\begin{equation}
\tau \sim \mathcal{U}(0,1),
\end{equation}
\begin{equation}
\mu_\tau = (1-\tau)z_{T-1} + \tau z_T,
\end{equation}
\begin{equation}
\sigma_\tau = \sigma_b \sqrt{\tau(1-\tau)},
\end{equation}
\begin{equation}
z_\tau = \mu_\tau + \sigma_\tau \epsilon, \qquad \epsilon \sim \mathcal{N}(0,I).
\end{equation}

The current progression flow forecaster is
\begin{equation}
(\hat{v}, \hat{z}_T, \widehat{\Delta t}, \hat{\sigma}, \hat{s})
=
F_\theta(z_\tau,\tau,c_h,c_p,c_s,z_{T-1}).
\end{equation}

The predicted future latent is parameterized as
\begin{equation}
\hat{z}_T = z_\tau + (1-\tau)\hat{v}.
\end{equation}

\subsection{Training Objectives}

The current total loss is
\begin{equation}
\mathcal{L} =
\lambda_v \mathcal{L}_{\mathrm{vel}}
+ \lambda_z \mathcal{L}_{\mathrm{target}}
+ \lambda_x \mathcal{L}_{\mathrm{recon}}
+ \lambda_t \mathcal{L}_{\mathrm{interval}}
+ \lambda_u \mathcal{L}_{\mathrm{uncertainty}}
+ \lambda_c \mathcal{L}_{\mathrm{consistency}}
+ \lambda_s \mathcal{L}_{\mathrm{structure}}.
\end{equation}

The individual terms are:
\begin{equation}
\mathcal{L}_{\mathrm{vel}} = \|\hat{v} - (z_T-z_{T-1})\|_2^2,
\end{equation}
\begin{equation}
\mathcal{L}_{\mathrm{target}} = \|\hat{z}_T - z_T\|_2^2,
\end{equation}
\begin{equation}
\mathcal{L}_{\mathrm{recon}} = \|\hat{x}_T - x_T\|_1,
\end{equation}
\begin{equation}
\mathcal{L}_{\mathrm{interval}} = (\widehat{\Delta t} - (t_T-t_{T-1}))^2,
\end{equation}
\begin{equation}
\mathcal{L}_{\mathrm{uncertainty}} = \frac{r}{\hat{\sigma}} + \log \hat{\sigma},
\end{equation}
where
\begin{equation}
r = \mathrm{mean}\big((\hat{z}_T-z_T)^2\big),
\end{equation}
\begin{equation}
\mathcal{L}_{\mathrm{consistency}} =
\|\Delta z_{\mathrm{pred}} - \Delta z_{\mathrm{hist}}\|_2^2,
\end{equation}
and
\begin{equation}
\mathcal{L}_{\mathrm{structure}} = \|\hat{s}-s^\star\|_2^2.
\end{equation}

\subsection{Inference}

At test time, the current implementation uses one-step forecasting from the previous latent:
\begin{equation}
\tau = 0, \qquad z_\tau = z_{T-1},
\end{equation}
followed by
\begin{equation}
(\hat{v}, \hat{z}_T, \widehat{\Delta t}, \hat{\sigma}, \hat{s})
=
F_\theta(z_{T-1},0,c_h,c_p,c_s,z_{T-1}).
\end{equation}

The model then decodes $\hat{z}_T$ to obtain the future image prediction.

\section{Experiments}

\subsection{Current Evaluation Scope}

The current implemented evaluation includes:
\begin{itemize}[leftmargin=1.5em]
  \item image quality metrics
  \item next-visit interval error
  \item uncertainty summary
  \item missing-visit robustness evaluation
\end{itemize}

\subsection{Placeholders for Result Tables}

\begin{table}[ht]
\centering
\caption{Placeholder for the main SCP-Flow results.}
\begin{tabular}{lcccc}
\toprule
Model & PSNR $\uparrow$ & SSIM $\uparrow$ & Interval MAE $\downarrow$ & FID $\downarrow$ \\
\midrule
SCP-Flow & -- & -- & -- & -- \\
\bottomrule
\end{tabular}
\end{table}

\begin{table}[ht]
\centering
\caption{Placeholder for ablation results of SCP-Flow components.}
\begin{tabular}{lccccc}
\toprule
Model Variant & Structure & Interval & Uncertainty & Consistency & Main Score \\
\midrule
SCP-Flow (base) & -- & -- & -- & -- & -- \\
SCP-Flow + structure & \checkmark & -- & -- & -- & -- \\
SCP-Flow + interval & -- & \checkmark & -- & -- & -- \\
SCP-Flow + uncertainty & -- & -- & \checkmark & -- & -- \\
SCP-Flow + consistency & -- & -- & -- & \checkmark & -- \\
\bottomrule
\end{tabular}
\end{table}

\begin{table}[ht]
\centering
\caption{Placeholder for missing-visit robustness results.}
\begin{tabular}{lcccc}
\toprule
Strategy & Count & PSNR $\uparrow$ & SSIM $\uparrow$ & Interval MAE $\downarrow$ \\
\midrule
none & 0 & -- & -- & -- \\
uniform & 1 & -- & -- & -- \\
random & 1 & -- & -- & -- \\
tail & 1 & -- & -- & -- \\
\bottomrule
\end{tabular}
\end{table}

\begin{figure}[ht]
\centering
\fbox{\rule{0pt}{1.8in}\rule{0.9\linewidth}{0pt}}
\caption{Placeholder for qualitative forecasting examples. Recommended layout: observed history, generated future image, ground-truth future image, and optionally structure ROI visualizations.}
\label{fig:qualitative}
\end{figure}


\section{Discussion}

The current SCP-Flow implementation is a minimum runnable prototype. Its core flow-based trajectory prediction, historical conditioning, structure conditioning, interval prediction, uncertainty head, and missing-visit evaluation are already connected end-to-end. However, several components remain prototype-level, especially uncertainty calibration and clinically grounded structure supervision.

\section{Conclusion}

This paper framework describes SCP-Flow, a structure-controllable progression flow model for longitudinal glaucoma forecasting. The current repository already supports the key implementation path of the proposed method, including latent flow prediction, historical and population conditioning, structure-aware inputs, interval prediction, uncertainty prediction, and missing-visit evaluation. The next step is to complete full-scale training and benchmark SCP-Flow against the original tHPM-LDM diffusion predictor under a unified evaluation protocol.


\bibliographystyle{plain}
\bibliography{references}

\end{document}
